\section{System Architecture}
\label{sec:architecture}

Fig.~\ref{fig:architecture} illustrates the end-to-end architecture of the
proposed NS-3 O-RAN 6G integration framework.

\begin{figure}[htbp]
  \centering
  %% TikZ architecture diagram – NS-3 O-RAN 6G Integration Framework
%% Included via %% TikZ architecture diagram – NS-3 O-RAN 6G Integration Framework
%% Included via %% TikZ architecture diagram – NS-3 O-RAN 6G Integration Framework
%% Included via \input{figures/architecture.tex} in main.tex
%% Run generate_paper_assets.py to regenerate figure stubs from results/.
\documentclass[tikz, border=15pt]{standalone}

%% Required Libraries
\usetikzlibrary{positioning, fit, arrows.meta, shadows, backgrounds, calc, shapes.multipart}

\begin{document}

\begin{tikzpicture}[
    font=\sffamily\small,
    >=Stealth,
    % ── Node Styles ──
    corebox/.style={
        draw=#1!80!black, thick, rounded corners=3pt,
        align=center, fill=#1!10, minimum height=0.8cm,
        drop shadow={opacity=0.1, shadow xshift=1pt, shadow yshift=-1pt}
    },
    xapp/.style={
        draw=violet!80!black, thick, rounded corners=2pt,
        align=center, fill=violet!5, minimum height=0.6cm, font=\scriptsize
    },
    % ── Boundary Styles ──
    domain/.style={
        draw=#1!60, thick, dashed, rounded corners=6pt, 
        inner sep=12pt, fill=#1!5, fill opacity=0.3, text opacity=1
    },
    % ── Interface Styles ──
    interface/.style={->, thick, draw=black!70, rounded corners=4pt},
    biinterface/.style={<->, thick, draw=black!70, rounded corners=4pt},
    lbl/.style={
        font=\sffamily\tiny\bfseries, midway, fill=white, 
        inner sep=1.5pt, draw=black!20, rounded corners=2pt
    },
    o1lbl/.style={
        font=\sffamily\tiny, near end, fill=white, 
        inner sep=1pt, text=black!60
    }
]

%% =========================================================
%% 1. SMO & Non-RT RIC (Cloud Layer)
%% =========================================================
\node[corebox=gray, minimum width=3.5cm, minimum height=1cm] (nonrt) at (6, 7) {\textbf{Non-RT RIC}\\ \scriptsize AI/ML Training \& Policy};
\node[corebox=gray, right=0.5cm of nonrt, minimum width=2.5cm, minimum height=1cm] (rapp) {\textbf{rApps}\\ \scriptsize Network DT \& Intent};

\begin{scope}[on background layer]
    \node[fit=(nonrt)(rapp), domain=gray, label={[font=\bfseries\small, text=gray!80!black]above:SMO (Service Management & Orchestration) - Cloud}] (smo_grp) {};
\end{scope}

%% =========================================================
%% 2. Near-RT RIC Domain (Edge Layer)
%% =========================================================
\node[corebox=red, minimum width=3.5cm, minimum height=3cm] (nearrt_base) at (8, 2.5) {};
\node[font=\bfseries, text=red!80!black, anchor=north] at (nearrt_base.north) {Near-RT RIC};

% xApps inside Near-RT RIC
\node[xapp, right=0.2cm of nearrt_base.west, yshift=0.3cm] (xapp_ho) {HO\\xApp};
\node[xapp, right=0.2cm of xapp_ho] (xapp_fl) {FL\\xApp};
\node[xapp, right=0.2cm of xapp_fl] (xapp_ts) {TS\\xApp};
\node[corebox=teal, minimum width=3.1cm, minimum height=0.6cm, above=0.2cm of nearrt_base.south, font=\scriptsize] (sdl) {Shared Data Layer (SDL) / DT};

\begin{scope}[on background layer]
    \node[fit=(nearrt_base), domain=red, inner sep=6pt] (ric_grp) {};
\end{scope}

%% =========================================================
%% 3. Multi-Access Edge Computing (MEC)
%% =========================================================
\node[corebox=orange, minimum width=2.5cm, minimum height=2.5cm] (mec) at (-1, 2.5) {};
\node[font=\bfseries, text=orange!80!black, anchor=north] at (mec.north) {MEC Host};
\node[corebox=white, minimum width=2cm, minimum height=0.6cm, font=\scriptsize, below=0.4cm of mec.north] (edge_ai) {Edge AI Server};
\node[corebox=white, minimum width=2cm, minimum height=0.6cm, font=\scriptsize, above=0.2cm of mec.south] (local_cache) {Local Cache};

\begin{scope}[on background layer]
    \node[fit=(mec), domain=orange, inner sep=6pt] (mec_grp) {};
\end{scope}

%% =========================================================
%% 4. Disaggregated 6G RAN (Cell Site)
%% =========================================================
% O-CU (Central Unit)
\node[corebox=green, minimum width=1.8cm] (ocu_cp) at (3.5, 3.2) {\textbf{O-CU-CP}};
\node[corebox=green, minimum width=1.8cm] (ocu_up) at (3.5, 1.8) {\textbf{O-CU-UP}};

% O-DU (Distributed Unit)
\node[corebox=green, minimum width=3.5cm] (odu) at (3.5, 0) {\textbf{O-DU}\\ \scriptsize High-PHY, MAC, RLC};

% O-RU (Radio Unit) & 6G Elements
\node[corebox=green, minimum width=2.5cm] (oru) at (3.5, -2) {\textbf{O-RU}\\ \scriptsize Low-PHY / Beamforming};
\node[corebox=cyan, minimum width=1.5cm, minimum height=0.6cm, font=\scriptsize, left=1cm of oru] (ris) {\textbf{RIS}\\ Controller};

\begin{scope}[on background layer]
    \node[fit=(ocu_cp)(ocu_up)(odu)(oru)(ris), domain=green, label={[font=\bfseries\small, text=green!80!black]above:Disaggregated 6G RAN}] (ran_grp) {};
\end{scope}

%% =========================================================
%% 5. User Equipment (UEs)
%% =========================================================
\node[corebox=blue, minimum width=1.5cm] (ue1) at (1.5, -4.5) {URLLC\\UE};
\node[corebox=blue, minimum width=1.5cm] (ue2) at (3.5, -4.5) {eMBB\\UE};
\node[corebox=blue, minimum width=1.5cm] (ue3) at (5.5, -4.5) {mMTC\\IoT};

\begin{scope}[on background layer]
    \node[fit=(ue1)(ue2)(ue3), domain=blue, label={[font=\bfseries\small, text=blue!80!black]below:User Equipment (UE) Swarm}] (ue_grp) {};
\end{scope}

%% =========================================================
%% 6. Standardized Interfaces (Routing)
%% =========================================================

% A1 Interface (Non-RT RIC <-> Near-RT RIC)
\draw[biinterface] (nonrt.south -| nearrt_base.north) -- node[lbl]{A1} (nearrt_base.north);

% O1 Interfaces (SMO to all network elements - dashed gray)
\draw[interface, dashed, draw=gray] (ocu_cp.west) -- ++(-0.5,0) |- node[o1lbl]{O1} (smo_grp.west);
\draw[interface, dashed, draw=gray] (ocu_up.west) -- ++(-0.7,0) |- (smo_grp.west);
\draw[interface, dashed, draw=gray] (odu.west)    -- ++(-0.9,0) |- (smo_grp.west);

% E2 Interfaces (RAN to Near-RT RIC)
\draw[interface] (ocu_cp.east) -- ++(1,0) |- node[lbl, near start]{E2} (nearrt_base.160);
\draw[interface] (ocu_up.east) -- ++(1.5,0) |- node[lbl, near start]{E2} (nearrt_base.180);
\draw[interface] (odu.east)    -- ++(2,0) |- node[lbl, near start]{E2} (nearrt_base.200);

% F1 Interfaces (CU to DU)
\draw[biinterface] (ocu_cp.south) -- node[lbl, left]{F1-c} (ocu_cp.south |- odu.north);
\draw[biinterface] (ocu_up.south) -- node[lbl, right]{F1-u} (ocu_up.south |- odu.north);

% Open Fronthaul (DU to RU)
\draw[biinterface] (odu.south) -- node[lbl]{Open Fronthaul (eCPRI)} (odu.south |- oru.north);

% RU to RIS
\draw[interface, dashed, draw=cyan!80!black] (oru.west) -- node[lbl, above]{Control} (ris.east);

% Air Interfaces (RU/RIS to UEs)
\draw[biinterface, draw=blue!70] (oru.south) -- node[lbl, sloped, above]{THz/Sub-THz} (ue2.north);
\draw[biinterface, draw=blue!70] (ris.south) -- node[lbl, sloped, above]{Reflected Path} (ue1.north);
\draw[biinterface, draw=blue!70] (oru.south) -- (ue3.north);

% MEC Routing (Data Plane)
\draw[biinterface] (mec.east |- ocu_up.west) -- node[lbl, above]{N6 / Local UPF} (ocu_up.west);

\end{tikzpicture}

\end{document}
 in main.tex
%% Run generate_paper_assets.py to regenerate figure stubs from results/.
\documentclass[tikz, border=15pt]{standalone}

%% Required Libraries
\usetikzlibrary{positioning, fit, arrows.meta, shadows, backgrounds, calc, shapes.multipart}

\begin{document}

\begin{tikzpicture}[
    font=\sffamily\small,
    >=Stealth,
    % ── Node Styles ──
    corebox/.style={
        draw=#1!80!black, thick, rounded corners=3pt,
        align=center, fill=#1!10, minimum height=0.8cm,
        drop shadow={opacity=0.1, shadow xshift=1pt, shadow yshift=-1pt}
    },
    xapp/.style={
        draw=violet!80!black, thick, rounded corners=2pt,
        align=center, fill=violet!5, minimum height=0.6cm, font=\scriptsize
    },
    % ── Boundary Styles ──
    domain/.style={
        draw=#1!60, thick, dashed, rounded corners=6pt, 
        inner sep=12pt, fill=#1!5, fill opacity=0.3, text opacity=1
    },
    % ── Interface Styles ──
    interface/.style={->, thick, draw=black!70, rounded corners=4pt},
    biinterface/.style={<->, thick, draw=black!70, rounded corners=4pt},
    lbl/.style={
        font=\sffamily\tiny\bfseries, midway, fill=white, 
        inner sep=1.5pt, draw=black!20, rounded corners=2pt
    },
    o1lbl/.style={
        font=\sffamily\tiny, near end, fill=white, 
        inner sep=1pt, text=black!60
    }
]

%% =========================================================
%% 1. SMO & Non-RT RIC (Cloud Layer)
%% =========================================================
\node[corebox=gray, minimum width=3.5cm, minimum height=1cm] (nonrt) at (6, 7) {\textbf{Non-RT RIC}\\ \scriptsize AI/ML Training \& Policy};
\node[corebox=gray, right=0.5cm of nonrt, minimum width=2.5cm, minimum height=1cm] (rapp) {\textbf{rApps}\\ \scriptsize Network DT \& Intent};

\begin{scope}[on background layer]
    \node[fit=(nonrt)(rapp), domain=gray, label={[font=\bfseries\small, text=gray!80!black]above:SMO (Service Management & Orchestration) - Cloud}] (smo_grp) {};
\end{scope}

%% =========================================================
%% 2. Near-RT RIC Domain (Edge Layer)
%% =========================================================
\node[corebox=red, minimum width=3.5cm, minimum height=3cm] (nearrt_base) at (8, 2.5) {};
\node[font=\bfseries, text=red!80!black, anchor=north] at (nearrt_base.north) {Near-RT RIC};

% xApps inside Near-RT RIC
\node[xapp, right=0.2cm of nearrt_base.west, yshift=0.3cm] (xapp_ho) {HO\\xApp};
\node[xapp, right=0.2cm of xapp_ho] (xapp_fl) {FL\\xApp};
\node[xapp, right=0.2cm of xapp_fl] (xapp_ts) {TS\\xApp};
\node[corebox=teal, minimum width=3.1cm, minimum height=0.6cm, above=0.2cm of nearrt_base.south, font=\scriptsize] (sdl) {Shared Data Layer (SDL) / DT};

\begin{scope}[on background layer]
    \node[fit=(nearrt_base), domain=red, inner sep=6pt] (ric_grp) {};
\end{scope}

%% =========================================================
%% 3. Multi-Access Edge Computing (MEC)
%% =========================================================
\node[corebox=orange, minimum width=2.5cm, minimum height=2.5cm] (mec) at (-1, 2.5) {};
\node[font=\bfseries, text=orange!80!black, anchor=north] at (mec.north) {MEC Host};
\node[corebox=white, minimum width=2cm, minimum height=0.6cm, font=\scriptsize, below=0.4cm of mec.north] (edge_ai) {Edge AI Server};
\node[corebox=white, minimum width=2cm, minimum height=0.6cm, font=\scriptsize, above=0.2cm of mec.south] (local_cache) {Local Cache};

\begin{scope}[on background layer]
    \node[fit=(mec), domain=orange, inner sep=6pt] (mec_grp) {};
\end{scope}

%% =========================================================
%% 4. Disaggregated 6G RAN (Cell Site)
%% =========================================================
% O-CU (Central Unit)
\node[corebox=green, minimum width=1.8cm] (ocu_cp) at (3.5, 3.2) {\textbf{O-CU-CP}};
\node[corebox=green, minimum width=1.8cm] (ocu_up) at (3.5, 1.8) {\textbf{O-CU-UP}};

% O-DU (Distributed Unit)
\node[corebox=green, minimum width=3.5cm] (odu) at (3.5, 0) {\textbf{O-DU}\\ \scriptsize High-PHY, MAC, RLC};

% O-RU (Radio Unit) & 6G Elements
\node[corebox=green, minimum width=2.5cm] (oru) at (3.5, -2) {\textbf{O-RU}\\ \scriptsize Low-PHY / Beamforming};
\node[corebox=cyan, minimum width=1.5cm, minimum height=0.6cm, font=\scriptsize, left=1cm of oru] (ris) {\textbf{RIS}\\ Controller};

\begin{scope}[on background layer]
    \node[fit=(ocu_cp)(ocu_up)(odu)(oru)(ris), domain=green, label={[font=\bfseries\small, text=green!80!black]above:Disaggregated 6G RAN}] (ran_grp) {};
\end{scope}

%% =========================================================
%% 5. User Equipment (UEs)
%% =========================================================
\node[corebox=blue, minimum width=1.5cm] (ue1) at (1.5, -4.5) {URLLC\\UE};
\node[corebox=blue, minimum width=1.5cm] (ue2) at (3.5, -4.5) {eMBB\\UE};
\node[corebox=blue, minimum width=1.5cm] (ue3) at (5.5, -4.5) {mMTC\\IoT};

\begin{scope}[on background layer]
    \node[fit=(ue1)(ue2)(ue3), domain=blue, label={[font=\bfseries\small, text=blue!80!black]below:User Equipment (UE) Swarm}] (ue_grp) {};
\end{scope}

%% =========================================================
%% 6. Standardized Interfaces (Routing)
%% =========================================================

% A1 Interface (Non-RT RIC <-> Near-RT RIC)
\draw[biinterface] (nonrt.south -| nearrt_base.north) -- node[lbl]{A1} (nearrt_base.north);

% O1 Interfaces (SMO to all network elements - dashed gray)
\draw[interface, dashed, draw=gray] (ocu_cp.west) -- ++(-0.5,0) |- node[o1lbl]{O1} (smo_grp.west);
\draw[interface, dashed, draw=gray] (ocu_up.west) -- ++(-0.7,0) |- (smo_grp.west);
\draw[interface, dashed, draw=gray] (odu.west)    -- ++(-0.9,0) |- (smo_grp.west);

% E2 Interfaces (RAN to Near-RT RIC)
\draw[interface] (ocu_cp.east) -- ++(1,0) |- node[lbl, near start]{E2} (nearrt_base.160);
\draw[interface] (ocu_up.east) -- ++(1.5,0) |- node[lbl, near start]{E2} (nearrt_base.180);
\draw[interface] (odu.east)    -- ++(2,0) |- node[lbl, near start]{E2} (nearrt_base.200);

% F1 Interfaces (CU to DU)
\draw[biinterface] (ocu_cp.south) -- node[lbl, left]{F1-c} (ocu_cp.south |- odu.north);
\draw[biinterface] (ocu_up.south) -- node[lbl, right]{F1-u} (ocu_up.south |- odu.north);

% Open Fronthaul (DU to RU)
\draw[biinterface] (odu.south) -- node[lbl]{Open Fronthaul (eCPRI)} (odu.south |- oru.north);

% RU to RIS
\draw[interface, dashed, draw=cyan!80!black] (oru.west) -- node[lbl, above]{Control} (ris.east);

% Air Interfaces (RU/RIS to UEs)
\draw[biinterface, draw=blue!70] (oru.south) -- node[lbl, sloped, above]{THz/Sub-THz} (ue2.north);
\draw[biinterface, draw=blue!70] (ris.south) -- node[lbl, sloped, above]{Reflected Path} (ue1.north);
\draw[biinterface, draw=blue!70] (oru.south) -- (ue3.north);

% MEC Routing (Data Plane)
\draw[biinterface] (mec.east |- ocu_up.west) -- node[lbl, above]{N6 / Local UPF} (ocu_up.west);

\end{tikzpicture}

\end{document}
 in main.tex
%% Run generate_paper_assets.py to regenerate figure stubs from results/.
\documentclass[tikz, border=15pt]{standalone}

%% Required Libraries
\usetikzlibrary{positioning, fit, arrows.meta, shadows, backgrounds, calc, shapes.multipart}

\begin{document}

\begin{tikzpicture}[
    font=\sffamily\small,
    >=Stealth,
    % ── Node Styles ──
    corebox/.style={
        draw=#1!80!black, thick, rounded corners=3pt,
        align=center, fill=#1!10, minimum height=0.8cm,
        drop shadow={opacity=0.1, shadow xshift=1pt, shadow yshift=-1pt}
    },
    xapp/.style={
        draw=violet!80!black, thick, rounded corners=2pt,
        align=center, fill=violet!5, minimum height=0.6cm, font=\scriptsize
    },
    % ── Boundary Styles ──
    domain/.style={
        draw=#1!60, thick, dashed, rounded corners=6pt, 
        inner sep=12pt, fill=#1!5, fill opacity=0.3, text opacity=1
    },
    % ── Interface Styles ──
    interface/.style={->, thick, draw=black!70, rounded corners=4pt},
    biinterface/.style={<->, thick, draw=black!70, rounded corners=4pt},
    lbl/.style={
        font=\sffamily\tiny\bfseries, midway, fill=white, 
        inner sep=1.5pt, draw=black!20, rounded corners=2pt
    },
    o1lbl/.style={
        font=\sffamily\tiny, near end, fill=white, 
        inner sep=1pt, text=black!60
    }
]

%% =========================================================
%% 1. SMO & Non-RT RIC (Cloud Layer)
%% =========================================================
\node[corebox=gray, minimum width=3.5cm, minimum height=1cm] (nonrt) at (6, 7) {\textbf{Non-RT RIC}\\ \scriptsize AI/ML Training \& Policy};
\node[corebox=gray, right=0.5cm of nonrt, minimum width=2.5cm, minimum height=1cm] (rapp) {\textbf{rApps}\\ \scriptsize Network DT \& Intent};

\begin{scope}[on background layer]
    \node[fit=(nonrt)(rapp), domain=gray, label={[font=\bfseries\small, text=gray!80!black]above:SMO (Service Management & Orchestration) - Cloud}] (smo_grp) {};
\end{scope}

%% =========================================================
%% 2. Near-RT RIC Domain (Edge Layer)
%% =========================================================
\node[corebox=red, minimum width=3.5cm, minimum height=3cm] (nearrt_base) at (8, 2.5) {};
\node[font=\bfseries, text=red!80!black, anchor=north] at (nearrt_base.north) {Near-RT RIC};

% xApps inside Near-RT RIC
\node[xapp, right=0.2cm of nearrt_base.west, yshift=0.3cm] (xapp_ho) {HO\\xApp};
\node[xapp, right=0.2cm of xapp_ho] (xapp_fl) {FL\\xApp};
\node[xapp, right=0.2cm of xapp_fl] (xapp_ts) {TS\\xApp};
\node[corebox=teal, minimum width=3.1cm, minimum height=0.6cm, above=0.2cm of nearrt_base.south, font=\scriptsize] (sdl) {Shared Data Layer (SDL) / DT};

\begin{scope}[on background layer]
    \node[fit=(nearrt_base), domain=red, inner sep=6pt] (ric_grp) {};
\end{scope}

%% =========================================================
%% 3. Multi-Access Edge Computing (MEC)
%% =========================================================
\node[corebox=orange, minimum width=2.5cm, minimum height=2.5cm] (mec) at (-1, 2.5) {};
\node[font=\bfseries, text=orange!80!black, anchor=north] at (mec.north) {MEC Host};
\node[corebox=white, minimum width=2cm, minimum height=0.6cm, font=\scriptsize, below=0.4cm of mec.north] (edge_ai) {Edge AI Server};
\node[corebox=white, minimum width=2cm, minimum height=0.6cm, font=\scriptsize, above=0.2cm of mec.south] (local_cache) {Local Cache};

\begin{scope}[on background layer]
    \node[fit=(mec), domain=orange, inner sep=6pt] (mec_grp) {};
\end{scope}

%% =========================================================
%% 4. Disaggregated 6G RAN (Cell Site)
%% =========================================================
% O-CU (Central Unit)
\node[corebox=green, minimum width=1.8cm] (ocu_cp) at (3.5, 3.2) {\textbf{O-CU-CP}};
\node[corebox=green, minimum width=1.8cm] (ocu_up) at (3.5, 1.8) {\textbf{O-CU-UP}};

% O-DU (Distributed Unit)
\node[corebox=green, minimum width=3.5cm] (odu) at (3.5, 0) {\textbf{O-DU}\\ \scriptsize High-PHY, MAC, RLC};

% O-RU (Radio Unit) & 6G Elements
\node[corebox=green, minimum width=2.5cm] (oru) at (3.5, -2) {\textbf{O-RU}\\ \scriptsize Low-PHY / Beamforming};
\node[corebox=cyan, minimum width=1.5cm, minimum height=0.6cm, font=\scriptsize, left=1cm of oru] (ris) {\textbf{RIS}\\ Controller};

\begin{scope}[on background layer]
    \node[fit=(ocu_cp)(ocu_up)(odu)(oru)(ris), domain=green, label={[font=\bfseries\small, text=green!80!black]above:Disaggregated 6G RAN}] (ran_grp) {};
\end{scope}

%% =========================================================
%% 5. User Equipment (UEs)
%% =========================================================
\node[corebox=blue, minimum width=1.5cm] (ue1) at (1.5, -4.5) {URLLC\\UE};
\node[corebox=blue, minimum width=1.5cm] (ue2) at (3.5, -4.5) {eMBB\\UE};
\node[corebox=blue, minimum width=1.5cm] (ue3) at (5.5, -4.5) {mMTC\\IoT};

\begin{scope}[on background layer]
    \node[fit=(ue1)(ue2)(ue3), domain=blue, label={[font=\bfseries\small, text=blue!80!black]below:User Equipment (UE) Swarm}] (ue_grp) {};
\end{scope}

%% =========================================================
%% 6. Standardized Interfaces (Routing)
%% =========================================================

% A1 Interface (Non-RT RIC <-> Near-RT RIC)
\draw[biinterface] (nonrt.south -| nearrt_base.north) -- node[lbl]{A1} (nearrt_base.north);

% O1 Interfaces (SMO to all network elements - dashed gray)
\draw[interface, dashed, draw=gray] (ocu_cp.west) -- ++(-0.5,0) |- node[o1lbl]{O1} (smo_grp.west);
\draw[interface, dashed, draw=gray] (ocu_up.west) -- ++(-0.7,0) |- (smo_grp.west);
\draw[interface, dashed, draw=gray] (odu.west)    -- ++(-0.9,0) |- (smo_grp.west);

% E2 Interfaces (RAN to Near-RT RIC)
\draw[interface] (ocu_cp.east) -- ++(1,0) |- node[lbl, near start]{E2} (nearrt_base.160);
\draw[interface] (ocu_up.east) -- ++(1.5,0) |- node[lbl, near start]{E2} (nearrt_base.180);
\draw[interface] (odu.east)    -- ++(2,0) |- node[lbl, near start]{E2} (nearrt_base.200);

% F1 Interfaces (CU to DU)
\draw[biinterface] (ocu_cp.south) -- node[lbl, left]{F1-c} (ocu_cp.south |- odu.north);
\draw[biinterface] (ocu_up.south) -- node[lbl, right]{F1-u} (ocu_up.south |- odu.north);

% Open Fronthaul (DU to RU)
\draw[biinterface] (odu.south) -- node[lbl]{Open Fronthaul (eCPRI)} (odu.south |- oru.north);

% RU to RIS
\draw[interface, dashed, draw=cyan!80!black] (oru.west) -- node[lbl, above]{Control} (ris.east);

% Air Interfaces (RU/RIS to UEs)
\draw[biinterface, draw=blue!70] (oru.south) -- node[lbl, sloped, above]{THz/Sub-THz} (ue2.north);
\draw[biinterface, draw=blue!70] (ris.south) -- node[lbl, sloped, above]{Reflected Path} (ue1.north);
\draw[biinterface, draw=blue!70] (oru.south) -- (ue3.north);

% MEC Routing (Data Plane)
\draw[biinterface] (mec.east |- ocu_up.west) -- node[lbl, above]{N6 / Local UPF} (ocu_up.west);

\end{tikzpicture}

\end{document}

  \caption{NS-3 O-RAN 6G integration framework architecture. UEs connect
           to gNBs over NR/THz air interfaces; gNBs report KPIs to the
           near-RT RIC via E2; xApps execute RL and FL control loops;
           the digital-twin module feeds back predicted KPIs.}
  \label{fig:architecture}
\end{figure}

\subsection{Radio Access Network Layer}
The RAN layer consists of \(N_{\text{gNB}}\) gNB nodes, each serving up
to \(N_{\text{UE}}\) user equipment (UE) nodes.
Channel models support sub-6\,GHz NR and THz bands (0.1--1\,THz).
Each gNB implements a Proportional Fair (PF) scheduler as the baseline and
an RL-enhanced scheduler as the advanced variant.

\subsection{Near-RT RIC and xApps}
The near-RT RIC is modelled as a centralised controller that receives
periodic E2 service model (SM) indication messages from all gNBs.
Two xApps are implemented:
\begin{itemize}
  \item \textbf{Handover xApp}: uses a Deep Q-Network (DQN) to decide
        source$\to$target handovers based on RSRP, RSRQ, and SINR reports.
  \item \textbf{Federated Learning xApp}: orchestrates FedAvg rounds across
        gNB-local models that predict UE throughput from radio measurements.
\end{itemize}

\subsection{Digital Twin Module}
A lightweight regression model (gradient-boosted tree) is trained online
using sliding-window measurements.
Before the RIC applies an action, the twin predicts the post-action KPI
and rejects actions predicted to worsen performance beyond a configurable
threshold.

\subsection{MEC Service Layer}
MEC hosts are co-located with gNBs (edge) and with a regional data centre
(cloud).
Services (Video Analytics, AR/VR Streaming, IoT Processing, FL model
aggregation) are modelled with measured CPU/memory utilisation and
variable latency depending on host tier and active users.

\subsection{Data Flow and Result Export}
Simulation outputs are written to \texttt{results/data/} as CSV and JSON
files.
The \texttt{tools/paper/generate\_paper\_assets.py} script reads these
outputs and produces \LaTeX{} tables and TikZ/PGFPlots figures that are
directly included in this document.
